\begin{figure*}[!t]
    \centering
    \includegraphics[width=\textwidth]{figures/assets_5_real_datasets_4_panels/appendix_reliability_weight_example.pdf}
    \caption{
    \textbf{Dataset-relative reliability weighting preserves
    observation-level variation.}
    \textbf{(A--D)}
    Each point represents an eligible transcript--dataset observation,
    positioned by positive-codon coverage $C_{dt}$ and
    log-transformed read density $\log(1+D_{dt})$.
    Background colors and contours indicate the normalized weight
    computed using the dataset's frozen training references
    $\tau_d$ and $m_d$.
    Examples were selected deterministically near the 12.5th,
    37.5th, 62.5th, and 87.5th percentiles of $\tau_d$ across
    the collection.
    \textbf{(E)}
    The dark-blue curve averages the 114 dataset-specific probability
    densities, assigning equal mass to each dataset regardless of
    transcript count.
    The blue envelope shows their pointwise 10th--90th percentiles,
    not a confidence interval; the pale vertical band marks the
    mixture interquartile range.
    Dashed red and dotted blue lines denote the normalization
    reference $w_{dt}=1$ and the mixture median, respectively;
    rug marks indicate dataset-specific medians.
    Equal dataset mass applies only to this descriptive visualization:
    training remains transcript-balanced.
    No model predictions or gamma-reference weights $\pi_d$
    enter the figure.
    }
    \label{fig:reliability_weight_example}
\end{figure*}

\subsection{Profile preprocessing and reliability weighting}
\label{sec:real_preprocessing}

Within each dataset, we retain nonempty, finite, nonnegative profiles
with positive total counts. All-zero profiles are excluded, while
sparse profiles and their zero-count codons are preserved.
Raw replicates remain separate for the NB2 loss.
For each transcript--dataset pair, the arithmetic replicate consensus
defines read density
$D_{dt}=n_t^{-1}\sum_i\overline Y_{dti}$
and positive-codon coverage
$C_{dt}=n_t^{-1}\sum_i\mathbf 1[\overline Y_{dti}>0]$,
where $n_t$ is the number of modeled codons.

The weighting is motivated by counting precision.
Under NB2, the mean-to-standard-deviation ratio is
$\mu/\sqrt{\mu+\alpha\mu^2}$: it grows approximately as
$\sqrt{\mu}$ at low counts and saturates at high counts for fixed
$\alpha>0$.
We therefore use a bounded square-root depth score, supplemented
by coverage to distinguish broad positional support from counts
concentrated at a few positions.
Let $\mathcal T_d^{\mathrm{train}}$ denote the retained training
transcripts in dataset $d$ and
$\tau_d=\operatorname{median}_{u\in\mathcal T_d^{\mathrm{train}}}D_{du}$.
We define
\begin{equation}
    \begin{aligned}
        w^{\mathrm{raw}}_{dt}
        &=
        0.70\,
        \frac{\sqrt{D_{dt}}}
             {\sqrt{D_{dt}}+\sqrt{\tau_d}}
        +0.30\,C_{dt},\\
        w_{dt}
        &=
        \frac{w^{\mathrm{raw}}_{dt}}{m_d},
        \qquad
        m_d=
        \operatorname{median}_{u\in\mathcal T_d^{\mathrm{train}}}
        w^{\mathrm{raw}}_{du}.
    \end{aligned}
    \label{eq:real_reliability}
\end{equation}
The depth component equals one-half at $D_{dt}=\tau_d$ and gives
diminishing returns to additional reads.
The mixture is an SNR-inspired reliability heuristic, not a
calibrated SNR or inverse-variance estimate.

Median normalization expresses each observation's reliability
relative to the typical training observation in its own dataset.
It preserves within-dataset differences while setting each dataset's
median training weight to one: weights above or below one indicate
above- or below-median sampling support within that dataset,
rather than absolute quality across experiments.
This removes differences in median raw scores as an additional
dataset-wide loss multiplier, but does not enforce equal total
dataset contributions.
Both $\tau_d$ and $m_d$ are fitted on training transcripts and frozen
for evaluation.
The resulting weights enter the transcript-balanced objective in
Equation~\ref{eq:transcript_balanced_objective}; they do not modify
the count profiles.
Their role is distinct from $\pi_d$, which defines the reference
priorities used to center the dataset-specific corrections.

Figure~\ref{fig:reliability_weight_example} illustrates this
dataset-relative interpretation.
Panels A--D show how the weighting rule adapts to datasets with
different reference densities, while panel E summarizes the
resulting distributions across all 114 datasets.
The displayed dataset medians lie between $0.996$ and $1.011$,
but substantial observation-level variation remains: the
equal-dataset mixture has an interquartile range of
$0.620$--$1.394$.
All eligible observations are displayed, including held-out rows
transformed without refitting the training references.
The near-unit medians are primarily a consequence of normalization,
not evidence of equal dataset quality or independent validation
of held-out reliability.

\subsection{Real-data stability experiments}
\label{sec:real_stability_experiments}

The real-data experiments study Ribo-seq collections with heterogeneous read
depth and unknown dataset-specific bias.  Without a known biological profile,
two claims remain identifiable.  First, independently fitted, source-disjoint
dataset panels should recover concordant $\mathbf{L}_t$ profiles for the
same held-out transcripts.  Second, a quality-ranked gamma reference should
reduce the displacement of $\mathbf{L}_t$ as progressively lower-ranked
datasets are added.  These are reproducibility and stability claims; real data
alone cannot establish recovery of the unknown true $\mathbf{L}_t$ or
$\gamma_{dt}$.

All experiments use 114 HEK Ribo-seq datasets aligned to the same MANE CDS
representation.  Replicates from one condition remain separate count profiles
within a transcript--dataset observation.  For dataset $d$, transcript $t$, and
codon $i$, the modeled mean is
\begin{equation}
    \mu_{dti}=S_{dt}\,[\mathbf{L}_t]_i\,\gamma_{dti},
    \qquad
    \frac{1}{n_t}\sum_i [\mathbf{L}_t]_i=1,
    \label{eq:real_mean_decomposition}
\end{equation}
where $S_{dt}$ is the target-derived abundance scale,
$\mathbf{L}_t$ is shared across datasets, and $\gamma_{dt}$ is the
positional dataset correction.

\subsubsection{Frozen dataset-quality ranking and gamma reference}
\label{sec:real_quality_ranking}

The dataset ordering is frozen before training.  It aggregates ten
prespecified QC indicators covering periodicity, CDS enrichment, depth and
transcript support, footprint-length quality, mapping and contamination, and
replicate agreement.  Rank $r_d=1$ denotes the best QC profile.  The ranking
table contains $R=115$ entries; 114 are active because global rank 110 is not
configured.  Ranks are never recomputed within a panel or cumulative subset.

For an active collection $\mathcal D$, the rank score and normalized reference
weight are
\begin{equation}
    q_d=\frac{R-r_d+1}{R},
    \qquad
    \pi_d^{(p)}=\frac{q_d^p}{\sum_{e\in\mathcal D}q_e^p}.
    \label{eq:real_gamma_reference}
\end{equation}
We use equal reference weights ($p=0$), ranked weights with $p=1$, and a
stronger $p=3$ sensitivity analysis.  Reversed controls attach the same
multiset of $q_d^p$ values to datasets in the opposite rank order.  They
therefore preserve weight concentration while reversing which datasets define
the reference.

The weights impose the gamma-identifiability constraint
\begin{equation}
    \sum_{d\in\mathcal D}\pi_d\log\gamma_{dti}=0
    \quad\text{for every }(t,i).
    \label{eq:real_gamma_centering}
\end{equation}
Thus $\pi_d$ changes how a fitted mean is divided between the shared profile and
the dataset-specific correction.  It does not remove low-ranked observations
or reweight their likelihood contributions; those contributions use the local
reliability weights $w_{dt}$ in Equation~\ref{eq:real_reliability}.

\subsubsection{Common training and comparison protocol}
\label{sec:real_stability_metrics}

Each model is initialized independently with seed 42 and selected by minimum
validation loss.  For transcript $t$, a batch group contains all usable
observations from the selected datasets; missing transcript--dataset pairs are
not fabricated.  Raw replicates enter the NB2 likelihood separately, whereas
their arithmetic consensus enters the profile-correlation terms.  Dataset
losses are combined within transcript using $w_{dt}$ and then averaged across
transcripts, so transcripts observed in more datasets do not automatically
receive greater outer weight.

The selected checkpoint exports a sequence-derived, full-CDS, mean-one
$\mathbf{L}_t$ profile.  After matching transcript IDs, CDS coordinates,
and valid-position masks, two models $A$ and $B$ are compared transcript by
transcript:
\begin{equation}
    \operatorname{PCC}_t(A,B)
    =\operatorname{corr}_{i}\!\left(
      [\mathbf{L}^{A}_t]_i,[\mathbf{L}^{B}_t]_i\right).
    \label{eq:real_profile_agreement}
\end{equation}
The distributions of these PCC values measure positional shape agreement.  We
also verify that the exported profiles retain within-transcript variation, so a
high PCC cannot be attributed to constant profiles.

\subsubsection{Four source-disjoint panels}
\label{sec:real_four_panel_design}

The first experiment partitions the 114 datasets once into four disjoint
panels.  Datasets inferred to share a publication or source family remain in
the same panel.  A seed-42, capacity-aware assignment balances quantile strata
of median read density, positive-codon coverage, eligible-transcript count, and
median replicate PCC, followed by size-preserving source-family swaps.  The
scalar QC rank is not used to construct the panels; it is used only to define
$\pi_d$ after membership is fixed.

\begin{table*}[!t]
    \centering
    \small
    \caption{Composition and training-cohort sizes of the four fixed
    source-disjoint panels.  All panels share 1,593 validation and 1,593 test
    transcript IDs.}
    \label{tab:real_four_panel_design}
    \begin{tabular}{lrrrrr}
        \toprule
        Panel & Datasets & Source families & Global-rank range & Mean rank & Training transcripts \\
        \midrule
        1 & 29 & 20 & 7--106 & 54.66 & 13,261 \\
        2 & 29 & 21 & 2--114 & 59.17 & 13,613 \\
        3 & 28 & 22 & 1--115 & 57.79 & 13,595 \\
        4 & 28 & 22 & 6--113 & 58.61 & 13,340 \\
        \bottomrule
    \end{tabular}
\end{table*}

A common evaluation pool contains transcripts with usable observations in at
least two datasets from every panel.  From 15,929 eligible transcripts, we draw
1,593 validation and 1,593 test IDs without replacement, stratified by
reliability decile and conserved-stalling-site category.  Both sets are removed
from every training fold.  Training eligibility then remains panel-specific,
as reported in Table~\ref{tab:real_four_panel_design}.  All five reference
policies reuse the same folds within a panel.

The comparison was audited at the exported-prediction level.  All 20 fits use
the same validation IDs for checkpoint selection and the same disjoint test IDs
for evaluation.  Every pairwise PCC uses identical test transcripts, CDS
coordinates, and masks.  The validation objectives still differ across panels
because their observed datasets differ, and the six panel pairs share four
fitted models; they are not six independent replications.

\begin{figure*}[!t]
    \centering
    \includegraphics[width=\textwidth]{figures/assets_5_real_datasets_4_panels/appendix_four_panel_rank_and_cross_panel_pcc.pdf}
    \caption{
    \textbf{Panel QC-rank balance and transcript-level concordance of the
    shared profile.}
    \textbf{(A)} Frozen global QC ranks in the four source-disjoint panels;
    points are datasets, boxes show the interquartile range and median, white
    diamonds show the mean, and the dashed line marks the active-collection
    mean.
    \textbf{(B--F)} PCC distributions for the same 1,593 test transcripts in
    each of the six panel pairs.  From left to right, the first row contains
    panel balance, ranked $p=1$, and ranked $p=3$; the second row contains equal
    reference, reversed $p=1$, and reversed $p=3$.  Embedded
    boxes show the median and interquartile range.  All PCC panels share
    adaptive limits computed from the pooled 0.5th and 99.5th percentiles with
    four-percent padding; the printed percentage gives the fraction outside
    those display limits, while all values remain in the summaries.
    }
    \label{fig:real_four_panel_overview}
\end{figure*}

Figure~\ref{fig:real_four_panel_overview} shows strong cross-panel concordance
under equal and mildly ranked centering, with no near-constant exported
profiles.  Agreement remains positive across all reference policies, although
strong $p=3$ centering---especially when reversed---reduces reproducibility.
The experiment therefore supports a shared, reproducible signal across
source-disjoint collections, while also showing that the latent decomposition
is not fully invariant to an aggressive gamma reference.  Because the panels
were balanced rather than ordered from best to worst, this experiment is not a
test that the QC ranking improves recovery.

\subsubsection{Controlled cumulative addition of lower-ranked datasets}
\label{sec:real_cumulative_design}

The cumulative experiment uses nested prefixes of the frozen ordering,
\begin{equation}
    N\in\{2,5,10,20,40,80,114\},
    \qquad
    \mathcal D_N=\{d_{(1)},\ldots,d_{(N)}\},
    \label{eq:real_cumulative_prefix}
\end{equation}
so each step retains the preceding datasets and adds the next lower-ranked
block.  Every size and reference policy is trained from a fresh initialization;
``cumulative'' refers only to membership.

Two deterministic diagnostics quantify how different the reference policies
actually are at each $N$:
\begin{equation}
    \frac{N_{\mathrm{eff}}}{N}
    =\frac{1}{N\sum_{d\in\mathcal D_N}\pi_d^2},
    \qquad
    \overline r_{\pi}=\sum_{d\in\mathcal D_N}\pi_d r_d.
    \label{eq:real_reference_diagnostics}
\end{equation}
The first measures concentration relative to equal weighting; the second shows
which end of the ranking receives that mass.  These quantities diagnose the
strength and direction of the intervention and are not model-performance
metrics.

\begin{figure*}[!t]
    \centering
    \includegraphics[width=\textwidth]{figures/assets_5_real_datasets_4_panels/appendix_cumulative_reference_concentration.pdf}
    \caption{
    \textbf{The cumulative ranking intervention becomes appreciable only for
    larger dataset collections.}
    \textbf{(A)} Effective reference fraction $N_{\mathrm{eff}}/N$ for each
    nested best-$N$ collection.  Ranked and reversed policies coincide because
    they contain the same weights.
    \textbf{(B)} Reference-weighted mean global rank, which separates their
    direction: ranked policies favor the best datasets and reversed policies
    favor the worst members of the same collection.  The shaded region marks
    $N\leq20$.  Both panels are fixed properties of the experimental design.
    }
    \label{fig:real_cumulative_reference_concentration}
\end{figure*}

Figure~\ref{fig:real_cumulative_reference_concentration} sets the expected
scale of the best-first result.  Through $N=20$, even the ranked references
within these top-ranked prefixes remain close to equal, so large early
differences in $\mathbf{L}_t$ would be surprising.  The intervention separates
mainly at the larger collections; those endpoints carry the strongest test of
whether reference direction stabilizes the shared profile.

The original best-first run evaluates every model on 1,771 common test
transcripts, although its training and validation cohorts vary with $N$.  A
second experiment fixes the cohort before training.  The intersection across
all 114 datasets yields 7,142 complete-case transcripts, split once into 5,714
training, 714 validation, and 714 test IDs.  Transcript membership, CDS
coordinates, observation availability, and training-fitted reliability
references are fixed across its collection sizes and policies.  This isolates
dataset addition, at the cost of restricting the analysis to transcripts
observed in every dataset.

\subsubsection{Bidirectional dataset-selection control}
\label{sec:real_selection_controls}

We combine the available best-first trajectory with the completed worst-first
branch of the fixed-cohort experiment.  The best-first path starts from the two
highest-ranked datasets and adds progressively lower-ranked datasets; the
worst-first path starts from the two lowest-ranked datasets and adds
progressively higher-ranked datasets.  Within each path, every model is
compared with one common equal-reference $N=2$ fit.  This shared anchor prevents
separately trained policy-specific anchors from entering the within-panel
ranking contrast.  It is a drift reference, not biological ground truth.

\begin{figure*}[!t]
    \centering
    \includegraphics[width=\textwidth]{figures/assets_5_real_datasets_4_panels/appendix_cumulative_ranking_direction_L_PCC.pdf}
    \caption{
    \textbf{Dataset-quality ranking moves the shared profile in the direction
    defined by the gamma reference.}
    \textbf{(A)} Best-first collections add progressively lower-ranked
    datasets.  Each point is the mean transcript-level PCC of
    $\mathbf{L}_t$ to the common equal-reference best-$N=2$ fit over the same
    1,771 test transcripts and CDS coordinates.  Ranked and reversed policies
    use the same selected datasets; at a fixed exponent they also have identical
    weight concentration and differ only in which end of the selected ranking
    receives that weight.  The shaded region marks $N\leq20$, where the ranked
    $p=3$ effective reference fraction remains at least $0.97$.  At $N=114$,
    equal and reversed $p=1$ were still unavailable and are not imputed.
    \textbf{(B)} Worst-first collections add progressively higher-ranked
    datasets.  PCC is measured to the common equal-reference worst-$N=2$ fit on
    the fixed cohort of 714 test transcripts.  Here $p=3$ is already
    concentrated at $N=2$ because the two lowest-ranked datasets have small,
    strongly separated quality scores.  In both panels, the equal $N=2$ point
    is the identity comparison; the other $N=2$ points compare their fitted
    policies with that same equal-reference anchor.  Bands are 95\% percentile
    intervals from a joint bootstrap over held-out transcripts within each
    panel and do not include optimization-seed uncertainty.  The two panels use
    adaptive vertical limits and different transcript cohorts, so their absolute
    slopes should not be compared quantitatively.
    }
    \label{fig:real_cumulative_ranking_direction}
\end{figure*}

Figure~\ref{fig:real_cumulative_ranking_direction} supports a directional
effect of the gamma reference.  Along the best-first path, ranked and reversed
$p=3$ remain close while the weights are nearly uniform.  After the collection
expands, ranked $p=3$ is consistently closer to the best-two anchor than its
concentration-matched reversed control, with the largest separation at
$N=114$.  The comparison isolates which datasets define the reference more
cleanly than ranked-versus-equal, whose difference remains modest and
non-monotonic.

The worst-first path provides the complementary check.  A quality-ranked
reference should not preserve a solution defined by the two worst datasets:
as better datasets enter, it assigns them greater reference mass and moves
$\mathbf{L}_t$ away from that poor anchor faster than equal weighting.  The
lower ranked-$p=3$ curve in panel B is therefore the expected directional
response, rather than evidence of worse stability.  Together, the two paths
show that $\pi_d$ controls where the model places the boundary between the
shared $\mathbf{L}_t$ profile and dataset-specific $\gamma_{dt}$ corrections in
a manner consistent with the QC ordering.

\subsubsection{Scope of the current evidence}

The completed four-panel experiment supports reproducibility of structured
shared profiles across source-disjoint real-data collections.  Conditional on
these fitted models, the cumulative results show a substantial,
direction-consistent effect of ranking on the latent decomposition; they do not
establish that the resulting $\mathbf{L}_t$ is closer to the unknown biological
truth.  That stronger claim requires known-signal recovery or an external
biological endpoint.  The best-first and worst-first panels also use different
cohorts and training protocols, and all runs use one optimization seed;
bootstrap intervals therefore do not measure optimization variability.

% Implementation provenance for this manuscript fragment:
% - four-panel design: results/four_panel_stability_seed42/experiment_manifest.json
% - controlled cumulative design: results/cumulative_stability_fixed_cohort_seed42/experiment_manifest.json
% - best/worst control: results/cumulative_selection_direction_seed42/experiment_manifest.json
% - quality-stratum control: results/four_quality_strata_seed42/experiment_manifest.json
% - four-panel result tables: analyses/artifacts/real_data/four_panel_stability/
% - exploratory cumulative tables: analyses/artifacts/real_data/cumulative_stability/
% - manuscript stability figures: analyses/create_iclr_real_data_stability_figures.py
% - cumulative direction figure: analyses/create_iclr_cumulative_ranking_direction_figure.py
% - frozen ranking SHA-256: 5811cadf68c56740e83b232b2990299d527630326205db9cba8bc7024d3cf1f8
